\section*{Brauhaus-Gulasch}
\ihead{}\chead{Personen:4}\ohead{}
\ifoot{Vorbereitungszeit:15 min}\cfoot{}\ofoot{Kochzeit:80 min}
{\Large Zutaten}
\begin{multicols}{2}
\begin{itemize}
\item \SI{500}{g} Schweinegulasch
\item \SI{2}{EL} Butterschmalz
\item \SI{1/2}{TL} Senf
\item \SI{1}{EL} Tomatenmark
\item \num{1} große Zwiebel
\item \num{1} Knoblauchzehe
\item \num{1} Karotte
\item \SI{200}{ml} Bier
\item \SI{600}{ml} Gemüsebrühe
\item Salz und Pfeffer
\item \num{1}{TL} Paprikapulver, edelsüß
\item \num{1}{EL} Crème fraîche
\item Saucenbinder nach Bedarf
\end{itemize}
% \columnbreak
\end{multicols}
\noindent {\Large Zubereitung}\ 
Das Butterschmalz in einer tiefen Pfanne erhitzen.
Das Gulasch im Butterschmalz scharf anbraten und mit Salz, Pfeffer und Paprikapulver würzen.
Wenn das Fleisch eine schöne Farbe angenommen hat und das Wasser verdampft ist, die gehackte Zwiebel, die gepresste Knoblauchzehe und die geriebene Karotte sowie den Senf und das Tomatenmark dazugeben und kurz weiterbraten.
Anschließend mit dem Bier ablöschen und einkochen lassen.
Mit der Gemüsebrühe auffüllen und mit Deckel ca. \SI{1}{h} köcheln lassen.
Etwa \SI{5}{min} vor Ende der Garzeit die Crème fraîche einrühren und nach Bedarf mit etwas Soßenbinder andicken.
Vor dem Servieren nochmal abschmecken.
Als Beilage passen Spätzle, Reis, Kartoffelbrei oder Salzkartoffeln.
